% !TEX root = ../main.tex

\chapter{系统需求分析}

\section{系统应用场景分析}

在实时数据平台中，Flink 作业通常承担持续运行的计算与写出职责，而 Kafka Topic 则经常随着业务边界变化而调整。
例如，多租户系统可能按租户自动创建 Topic，日志平台可能按日期、地域或业务线拆分 Topic，
灰度迁移和容灾切换场景中也可能同时存在新旧 Topic 或新旧集群。若仍采用传统静态 Sink，
作业在提交时就把目标 Topic 与连接参数固定下来，一旦下游目标变化，就需要停止作业、修改配置并重新部署。
这一过程会引入运维窗口、恢复等待时间和潜在重复写入风险，与实时链路“持续写入、低人工介入、快速接纳新目标”的要求相矛盾。

因此，本文面对的核心问题不是“如何把数据写入一个 Kafka Topic”，而是“当可写 Topic 集合在运行期间变化时，
Flink Sink 如何在不重启作业的前提下发现变化、更新写入目标，并保持状态与提交路径可恢复”。
这一问题天然包含两个层次：第一，系统需要从 Kafka 集群元数据中识别符合规则的新 Topic；
第二，系统需要把业务侧的逻辑路由与实际物理 Topic 解耦，使业务记录不必直接绑定某个固定目标。
前者对应自动发现路径，后者对应显式逻辑路由路径。

在本项目中，connector 通过 Topic 匹配模式、路由映射和元数据发现间隔等参数暴露动态能力，
而应用层只负责把这些参数装配进运行中的 Flink 作业。
因此，系统不仅要支持“自动发现符合模式的 Topic”，还要支持“根据业务路由标识动态决定消息应写往哪一类 Topic”。
这意味着问题已从“静态写一个 Topic”扩展为“运行期维护逻辑路由、物理 Topic 与底层写入通道的关系”。

与抽象需求相对应，当前仓库中的配置已经给出了较为完整的输入形式，本文将其整理至附录~\ref{appendix:configs}。
从这些配置可以看出，该 connector 同时暴露了自动发现所需的 \id{stream_pattern} 与
\id{discovery_interval_ms}，以及显式路由所需的 \id{route_map}。这意味着本文讨论的并不是
单一路径下的动态写入，而是“自动发现路径”和“逻辑路由路径”并存的复合场景。

\section{功能需求分析}

由上述场景可知，静态 Sink 的不足并不只是一项配置不够灵活，而是会沿着“目标发现、目标选择、写入通道、状态恢复”
形成连续影响：如果系统无法发现新 Topic，就不能减少停机变更；如果无法把逻辑路由映射到物理 Topic，
业务记录仍会被固定目标约束；如果发现目标后不能自动维护写入通道，数据仍无法真正写出；如果路由状态不能进入
Checkpoint，故障恢复后又会丢失运行期动态信息。因此，本节将功能需求拆分为动态 Topic 发现、自动订阅、
动态路由和数据处理四类，它们分别对应前述问题链条中的不同环节，并共同支撑后续第 4 章的系统设计。

\subsection{动态 Topic 发现需求}

系统应能根据用户配置的 Topic 名称模式，周期性查询 Kafka 当前 Topic 列表，并识别所有符合规则的目标。
当前实现中，构建阶段接收用户配置的 Topic 匹配模式，运行期由订阅规则调用元数据服务完成匹配；
在单集群原型中，内部流标识直接等价于 Topic 名称，因此模式匹配实质上就是对集群 Topic 名称做运行期过滤。

\subsection{自动订阅需求}

Sink 侧的“订阅”不是消费者组意义上的分区订阅，而是指写入执行模块对“当前可写目标集合”的维护。
在实现中，这一集合由内部活动路由集合表示；当元数据刷新触发时，系统重新计算新的可写路由集合，再按路由标识检查是否需要创建新的底层写入通道。
因此，本系统必须支持以下功能：一是首次启动时根据元数据建立可写集合；二是在运行期发现新增路由时自动扩展写入路径；
三是在路由不再被活动集合或显式逻辑路由引用时，尽可能释放对应写入通道，避免长时间累积陈旧资源。

需要指出的是，当前代码已经支持路由的新增与切换，并在无事务和至少一次路径下补充了退役写入通道清理；
但在精准一次路径下，由于事务上下文可能跨 Checkpoint 存在，当前实现仍采用保守策略，不自动清理退役写入通道。
因此，自动订阅需求在本文原型中已经形成基本闭环，但事务路径下的安全回收仍属于后续工程化完善方向。

\subsection{动态路由需求}

系统应支持“逻辑路由标识 $\to$ 物理 Kafka 目标”的两级映射，而不是把业务记录直接绑定到固定 Topic。
在应用层，输入事件同时覆盖数据记录与路由更新两种语义：普通数据事件携带路由标识与消息内容，更新事件携带路由到目标描述的映射。
在运行时，写入执行模块会先判断输入事件属于控制面还是数据面；若是更新事件则更新逻辑路由映射；若是普通事件且显式指定了路由标识，则走逻辑路由路径。

这说明本系统的路由需求具有两个层面：其一，面向“自动发现”的通用路由；其二，面向“显式路由标识”
的逻辑路由。后者正是本文“逻辑路由与物理路由分离”设计思想的直接代码体现。

进一步地，当前原型并没有把路由配置静态地塞进 Sink，而是通过一个轻量级验证载体利用广播状态与控制事件把路由更新传播到运行中的算子实例。
图~\ref{fig:route-broadcast-flow} 展示了这一机制的关键处理逻辑：普通数据事件只有在广播状态中存在对应路由标识时才会被放行，而路由更新事件则会先改写广播状态，再继续向下游 Sink 传播。
需要强调的是，这一部分并非 connector 的核心交付物，而是为了验证 connector 的显式路由更新接口而构造的上层驱动。

\begin{figure}[htbp]
  \centering
  \begin{tikzpicture}[
    node distance=0.85cm and 1.2cm,
    box/.style={draw, rounded corners=2pt, align=center, minimum width=2.7cm, minimum height=0.8cm},
    arrow/.style={-Latex, thick}
  ]
    \node[box] (update) {路由更新事件};
    \node[box, right=of update] (write) {写入广播状态\\更新路由配置};
    \node[box, right=of write] (forward) {继续转发\\给 Sink};
    \node[box, below=of update] (data) {普通数据事件\\携带路由标识};
    \node[box, right=of data] (check) {检查广播状态\\是否存在路由};
    \node[box, right=of check] (pass) {存在则\\放行数据进入 Sink};
    \node[box, below=of check] (drop) {不存在则过滤};
    \draw[arrow] (update) -- (write);
    \draw[arrow] (write) -- (forward);
    \draw[arrow] (data) -- (check);
    \draw[arrow] (check) -- node[above, font=\small] {存在} (pass);
    \draw[arrow] (check) -- (drop);
  \end{tikzpicture}
  \caption{广播状态维护路由配置与转发流程}
  \label{fig:route-broadcast-flow}
\end{figure}

这一实现意味着系统的功能需求实际上还隐含了一个重要约束：控制面配置必须先于数据面生效，
否则普通数据即使携带了合法的逻辑路由标识，也会因为广播状态中尚不存在该路由而被过滤掉。
因此，本文在分析动态路由需求时，需要同时考虑“路由如何表示”和“路由如何传播”两个问题。

\subsection{数据处理需求}

系统应能在持续数据流上稳定完成以下处理链路：数据生成/采集、路由校验、路由解析、Topic 选择、
消息序列化、Kafka 写入、Checkpoint 快照与故障恢复。当前验证载体使用数据生成源产生测试数据，并通过广播处理逻辑先过滤掉不存在的路由标识，再将合规事件送入动态 Sink。
说明系统在验证层面已经具备“数据面”和“控制面”混合流的基本处理能力，而 connector 本身则只关心进入写入模块之后的解析、写入与状态管理。

从实现要求看，序列化层还必须适应动态 Topic 绑定。当前实现会为每个路由维护对应的序列化与目标 Topic 绑定关系，
使普通消息在保留原有 key、value 与 headers 的同时，被写入解析后的目标 Topic。
这一机制是当前原型可运行的关键约束之一。

\section{非功能需求分析}

为了避免非功能需求停留在定性描述，本文将其限定为当前原型能够在本地实验中验证的指标。
这些指标不是生产级服务等级协议，而是用于判断本文方案是否满足毕业设计原型目标的可测约束。
各项指标的取值依据来自三类因素：其一是第 3.1 节场景分析给出的业务约束，即实时链路希望在不重启作业的情况下
快速接纳新增 Topic，且动态能力不能显著破坏持续写入能力；其二是当前实现的机制上界，例如
\id{discovery_interval_ms=2000} 决定了新增 Topic 通常需要等待下一轮刷新后才能被发现；
其三是本地单机原型的实验边界，例如第 5 章采用 \id{parallelism=1}、小规模路由数量和约 30 秒运行窗口进行验证。
因此，本节中的数值不是孤立给出的，而是由“场景诉求—实现机制—实验条件”共同推导得到。

\subsection{非功能指标的推导}

首先，发现实时性指标来源于“业务目标变化后不应再以停机重启为主要生效方式”的场景诉求。
在传统静态 Sink 中，新增 Topic 往往需要修改配置并重新提交作业，生效时间通常以人工操作和作业恢复时间计。
本文原型的目标不是把发现延迟压到毫秒级，而是证明在默认轮询机制下，Topic 变化可以在一个很短的自动化窗口内被接纳。
由于实验配置采用 \id{discovery_interval_ms=2000}，理论上新增 Topic 最坏会等待接近一个刷新周期；
再考虑 Kafka 管理面查询、写入通道创建和本地任务调度开销，本文额外预留约 \id{1 s}，
因此将“首次写入不超过 \id{3 s}”作为发现实时性目标。该阈值既来自当前实现的刷新周期，又明显小于停机重启的运维窗口，
能够体现动态发现机制的实际价值。

其次，吞吐可接受性和小规模扩展性指标来源于“动态能力不应以明显牺牲稳定写入能力为代价”的场景诉求。
本文方案相比静态 KafkaSink 增加了路由查找、正则匹配结果缓存、多写入通道管理和提交对象包装，
因此完全没有开销并不现实；但在本地单机、两路由到四路由的原型规模下，这些逻辑应当只是附加管理成本，不应成为主要吞吐瓶颈。
基于这一判断，本文将动态两路由相对静态基线的吞吐下降阈值设为 5\%，用来判断动态管理层是否仍处于可接受开销范围；
将两路由增至四路由的下降阈值设为 10\%，用来观察小规模路由扩展是否立即导致写入路径退化。
若超过这些阈值，则说明路由管理、写入通道创建或提交包装已经对原型目标产生明显影响，需要在设计层面重新考虑缓存、批量刷新和写入通道复用策略。

再次，Checkpoint 协同指标来源于“动态路由状态必须进入 Flink 原生容错链路”的场景诉求。
本文并不要求本地原型达到生产集群的大规模容错性能，但要求路由级状态快照不能明显阻塞小规模作业。
在实验中 Checkpoint 间隔设置为 \id{5000 ms}，因此稳态单次 Checkpoint 若达到数百毫秒甚至秒级，
就会说明动态状态封装已经对正常流处理产生可见干扰。本文将稳态完成时延阈值设为 \id{100 ms}，
是为了保证快照耗时相对 5 秒间隔保持在较低比例内，并便于通过第 5 章日志直接观察。

最后，事务路径可运行性指标来源于“动态 Sink 不能只在至少一次语义下成立”的场景诉求。
完整精准一次语义需要故障注入、恢复验证和下游 \id{read_committed} 消费检查共同支撑，
这超出了本文本地原型实验的覆盖范围。因此本文没有把该指标定义为生产级端到端精准一次证明，
而是定义为更保守的可运行性目标：精准一次小规模运行中至少完成 5 次 Checkpoint，
并在两个目标 Topic 上观察到 committed 输出。这个指标用于证明路由级事务前缀、提交对象包装和提交分发链路能够跑通，
为后续更严格故障验证奠定基础。表~\ref{tab:nfr-targets} 汇总了上述量化指标及其验证位置。

\begin{table}[htbp]
  \centering
  \caption{本文原型的非功能需求指标}
  \label{tab:nfr-targets}
  \begin{tabular}{p{0.18\textwidth}p{0.34\textwidth}p{0.38\textwidth}}
    \toprule
    需求类型 & 量化指标 & 指标依据与验证方式 \\
    \midrule
    发现实时性 & 在发现间隔为 2000ms 的条件下，新增匹配 Topic 的首次写入延迟应不超过 \id{3 s}。 &
    理论上延迟主要由下一轮刷新等待时间、Kafka 管理面查询和任务调度组成；第 5 章通过发现实时性实验验证。 \\
    吞吐可接受性 & 在本地单机、至少一次、两路由显式路由场景下，动态方案吞吐较静态基线下降不超过 5\%。 &
    动态 Sink 增加路由解析和写入通道查找，但不应在小规模路由下成为主导瓶颈；第 5 章静态基线对照实验验证。 \\
    小规模扩展性 & 在相同本地环境下，路由数量由 2 增至 4 时，吞吐不应出现超过 10\% 的下降。 &
    该指标用于验证小规模多路由管理不会立刻压垮写入路径；第 5 章吞吐测试表进行验证。 \\
    Checkpoint 协同 & 打开间隔为 \id{5000 ms} 的 Checkpoint 后，作业应能连续完成 Checkpoint，稳态完成时延控制在 \id{100 ms} 以内。 &
    路由级状态快照不应明显阻塞本地小规模作业；第 5 章 Checkpoint 时延实验验证。 \\
    事务路径可运行性 & 精准一次小规模运行中应完成至少 5 次 Checkpoint，并在两个目标 Topic 均产生 committed 输出。 &
    该指标验证路由级事务提交链路可运行，但不等同于复杂故障注入下的完整精准一次证明；第 5 章 Exactly-once 验证对应。 \\
    \bottomrule
  \end{tabular}
\end{table}

\subsection{实时性要求}

实时性的核心不是单条消息发送延迟，而是“Topic 变化被发现并投入写入”的时间上界。
当前实现通过配置项 \id{stream-metadata-discovery-interval-ms} 控制元数据刷新周期；
该参数默认可设置为非正值以关闭自动发现，应用层构建 Sink 时会把动态发现间隔配置传入该参数，示例配置默认是 \id{2000 ms}。
因此，系统的最小发现粒度大致受该周期限制，再叠加 Kafka 管理面查询返回时间与作业线程调度延迟，形成端到端发现时延。
基于这一机制，本文将本地原型的实时性目标定义为：默认 \id{2000 ms} 刷新周期下，
新增匹配 Topic 从创建成功到首次出现写入的时间不超过 \id{3 s}。该阈值比配置周期多预留约 \id{1 s}
给 Kafka 管理面调用与本地任务调度，既符合当前轮询实现的理论上界，也能在第 5 章通过 offset 观测直接验证。

\subsection{可扩展性要求}

可扩展性主要体现在三个方面。首先，随着匹配 Topic 数量增加，内部维护的底层写入通道数量同步增加，Kafka 连接、事务上下文和缓存占用也会提高。
其次，路由解析需要对已订阅流和集群元数据做展开，因此当 Topic 与集群数量扩大时，刷新计算本身会带来额外成本。
再次，路由标识目前由集群标识、Topic 与连接地址等信息拼接生成，
虽然实现简单、可唯一标识物理目标，但在大规模场景下需要更系统的命名与管理策略。

考虑到本文实验环境是本地单机原型，而非生产集群压测，本文把可扩展性指标限定在小规模路由扩展是否引入明显退化：
一是在两路由显式路由场景下，动态方案较静态基线吞吐下降不超过 5\%；
二是在相同条件下将路由数量从 2 增加到 4 时，吞吐下降不超过 10\%。
这两个阈值的合理性在于：动态 Sink 相比静态 Sink 必然增加路由查找、写入通道管理和提交包装开销，
但在毕业设计原型所覆盖的小规模场景中，这些开销不应主导整体吞吐；若下降超过上述范围，
则说明动态管理逻辑已经成为主要瓶颈，需要在设计阶段重新考虑缓存、批量刷新或写入通道复用策略。

\subsection{稳定性与容错性要求}

动态写入系统不能只关注“能否写出去”，还必须保证故障后状态可恢复。当前实现中，动态 Sink 接入 Flink 统一 Sink API 的状态化两阶段提交路径。
快照阶段会输出路由级恢复状态，其中包含路由标识、目标集群、Topic、事务前缀、producer 配置以及底层写入状态；
提交阶段则生成带路由元数据的提交对象，并按路由分发给底层提交路径。
因此，系统至少需要满足：快照内容足以重建路由级写入通道，恢复后能够重新拉起通道并继续工作，提交路径能够按路由维度保持隔离。

本文对稳定性与容错性的原型级量化要求包括两项。第一，在开启间隔为 \id{5000 ms} 的 Checkpoint 后，
作业应能够连续完成 Checkpoint，且排除首次冷启动影响后的稳态 Checkpoint 完成时延应控制在 \id{100 ms} 以内；
该指标用于判断路由级状态封装是否给本地小规模作业引入明显阻塞。第二，在
精准一次小规模验证中，系统应至少完成 5 次 Checkpoint，并在两个目标 Topic 上观察到 committed 输出；
该指标用于证明路由级事务提交路径能够跑通。需要强调的是，后者并不等同于覆盖网络中断、TaskManager 崩溃、
事务超时等复杂故障的完整精准一次证明，而是本文原型阶段可验证的最低一致性闭环。
与此同时，日志、监控指标、异常处理与精准一次路径下退役写入通道的安全清理等工程能力，
仍是进一步增强稳定性的重点。

\section{问题建模与分析}

设运行时 Kafka Topic 集合为 $T(t)$，由正则模式 $P$ 过滤出的自动发现集合为
$S(t)=\{\tau\in T(t)\mid \tau\ \text{匹配}\ P\}$。对于显式路由模式，系统还维护逻辑路由集合
$R=\{r_1,r_2,\dots,r_n\}$ 以及映射函数
$g:R\rightarrow 2^{S(t)}$，表示一个逻辑路由标识在时刻 $t$ 可以映射到的物理 Topic 集合。
对任意输入记录 $e$，若其未显式指定路由标识，则走自动发现路径并广播写入当前活动路由集合；
若其显式指定逻辑路由标识，则写入集合由 $g(r,t)$ 决定。

当前原型的一个重要实现特征是“自动发现路径”和“显式逻辑路由路径”并存：
前者基于订阅规则和元数据自动展开，后者基于路由目标描述中的连接地址与 Topic 模式解析得到。
因此，系统本质上是在同一写入执行模块内维护两类路由解析逻辑，并保证两者都能落入统一的快照与提交框架中。

\clearpage
\section{本章小结}

本章从实时链路长期运行与 Kafka Topic 动态变化之间的矛盾出发，独立推导了本文系统需要解决的问题：
既要在运行期发现符合规则的新 Topic，也要将业务逻辑路由与物理 Topic 解耦，并让写入通道、状态和提交路径跟随路由
变化而调整。随后，本章将这些问题拆解为动态 Topic 发现、自动订阅、动态路由和数据处理四类功能需求，
并进一步给出发现实时性、吞吐可接受性、小规模扩展性、Checkpoint 协同和事务路径可运行性等量化非功能指标。
这些指标将在第 4 章的设计中分别对应到元数据刷新、路由解析、写入通道管理、状态快照和提交协调设计，
并在第 5 章通过发现实时性、吞吐可接受性、小规模扩展性、Checkpoint 协同与事务路径可运行性实验进行验证，从而形成需求、设计与实验之间的闭环。
